% ============================================================
%  Introduction Générale
% ============================================================

\chapter*{Introduction}
\addcontentsline{toc}{chapter}{Introduction}
\markboth{Introduction}{}

La cartographie des cultures agricoles à grande échelle est un enjeu central
pour la gestion des ressources alimentaires et le suivi environnemental. Les
séries temporelles d'images satellitaires multibandes — notamment Sentinel-2
avec ses 10 bandes spectrales et ses revisites décadaires — offrent une source
d'information riche pour discriminer les cultures par leur signature phénologique.
Cependant, la variabilité inter-régionale, les données manquantes liées à la
couverture nuageuse, et la diversité des architectures de cultures rendent la
tâche de classification difficile.

Ce projet s'inscrit dans le cadre du module \textit{Réseaux de Neurones} et
porte sur la classification de cultures agricoles à partir de séries temporelles
d'images satellitaires Sentinel-2. Le point de départ est la reproduction et
l'analyse du modèle MCTNet (\textit{Multi-stage CNN-Transformer Network})
proposé par Wang et al.~\cite{wang2024}, qui combine des blocs CNN et Transformer
en parallèle pour capturer conjointement les dépendances locales et globales
des séries temporelles.

Le travail est organisé en trois parties complémentaires :

\begin{itemize}
  \item \textbf{Partie 1} — Reproduction complète de MCTNet : collecte des
        données via Google Earth Engine sur deux régions agricoles contrastées
        (Arkansas, 5 classes ; Californie, 6 classes), exploration et
        prétraitement, entraînement et évaluation. Les performances obtenues
        (OA\,=\,0.9603 Arkansas, OA\,=\,0.9311 Californie) sont cohérentes
        avec les résultats publiés.

  \item \textbf{Partie 2} — Intégration de covariables environnementales :
        enrichissement du modèle MCTNet par trois groupes de données exogènes
        (propriétés pédologiques OpenLandMap, variables climatiques ERA5,
        topographie ETOPO1/ALOS) et étude d'ablation systématique sur 7
        configurations.

  \item \textbf{Partie 3} — Trois contributions architecturales indépendantes :
        \textbf{GatedMCTNet} (fusion CNN-Transformer dynamique par gate apprise),
        \textbf{MCTNet Multiscale} (branche CNN multi-échelle à trois champs
        réceptifs), et \textbf{MCTNetUSkip} (connexions skip préservant les
        représentations à haute résolution temporelle, avec extension à un
        encodeur-décodeur U-Net et covariables).
\end{itemize}

\cleardoublepage
